\section{Lo simple es bello: retirar las capas intermedias y también los niveles jerárquicos de la organización}

Lo anterior trató la eliminación en el plano técnico; esta sección pasa a la eliminación en el plano organizacional. El «lo simple es bello» de Liu Xianming (刘先明) no es un lema estético, sino un principio de gestión de ingeniería: para que un sistema complejo itere con rapidez, es necesario reducir las capas intermedias innecesarias.

Esta sección pasa del «retirar» técnico a la «simplificación» organizacional. Liu Xianming dice que le gusta lo simple y que, al asumir el cargo, su decisión más importante fue simplificar los procesos, simplificar las etapas de investigación y desarrollo, fusionar las tareas duplicadas y bajar la prioridad de lo innecesario. No le gusta limitarse a escuchar informes; prefiere revisar directamente los problemas de primera línea, el código y los resultados de los experimentos.

\begin{figure}[H]
\centering
\includegraphics[width=0.96\textwidth]{simple-is-beautiful.png}
\caption{Lo simple es bello: retirar una y otra vez las densas capas intermedias para que los datos entrenen directamente las capacidades. Diagrama conceptual propio, elaborado a partir de la conversación de 00:00:00--00:02:16 y 00:54:30--01:48:46.}
\end{figure}

\begin{knowledgebox}{Lectura de la figura: la simplificación técnica y la simplificación organizacional son lo mismo}
Retirar la dependencia del LiDAR, retirar las reglas de planificación y control, retirar las capas intermedias y retirar Language tienen como fin reducir cuellos de botella; en la organización, reducir niveles jerárquicos, revisar directamente los problemas y fusionar el trabajo duplicado también buscan reducir los cuellos de botella de información.
\end{knowledgebox}

\subsection{Organización de ingeniería plana}

En escenarios donde la IA física se combina con la producción en serie, la velocidad de decisión es importante. Liu Xianming busca aplanar la jerarquía: que los engineer de primera línea puedan exponer los problemas directamente y que los lead de equipo también puedan escribir código, revisar experimentos y hacer deep dive. Dentro de la organización no se levantan muchos muros entre departamentos, para evitar el duplicate work, y la velocidad global aumenta al compartir experimentos y fracasos.

\begin{figure}[H]
\centering
\includegraphics[width=0.96\textwidth]{org-flat-engineering.png}
\caption{Organización de ingeniería plana: combinar la IA de frontera con la producción en serie exige decidir con rapidez sobre los problemas de primera línea. Diagrama conceptual propio, elaborado a partir de la conversación de 00:54:30--01:48:46.}
\end{figure}

\begin{knowledgebox}{Lectura de la figura: plano no significa sin gestión, sino cadena corta}
Los ingenieros de primera línea exponen código y experimentos, los responsables revisan directamente los problemas, el deep dive analiza los experimentos buenos y malos, la integración de recursos reduce la duplicación y, al final, todo ello permite un cambio rápido una vez que el modelo supera a las reglas.
\end{knowledgebox}

\subsection{Cultura de ingeniería: Deep Dive, compartir los fracasos y reducir la duplicación}

Esta subsección convierte lo «plano» en mecanismos cotidianos. Sin deep dive ni intercambio de fracasos, una organización plana corre el riesgo de tener solo unos cuantos niveles de reporte menos, en lugar de aumentar de verdad la velocidad de aprendizaje.

En la entrevista se menciona que el equipo hace deep dive: se comparten los experimentos buenos, los experimentos malos, lo que se ha visto recientemente y lo que a cada quien le salió mal. Esta práctica sirve a dos objetivos: primero, que el equipo forme rápidamente un contexto común; segundo, evitar que distintos grupos reinventen la rueda. La investigación y el desarrollo de Physical AI requieren una gran cantidad de experimentos; si los fracasos no se comparten, la organización cae una y otra vez en el mismo error.

\begin{importantbox}{Reutilización del conocimiento en la organización de ingeniería}
El activo central de un equipo de IA de frontera no son solo los pesos del modelo, sino también los experimentos fallidos, la experiencia en el ajuste de hiperparámetros, los métodos de procesamiento de datos y la comprensión compartida de los problemas. El Deep Dive vuelve explícito este conocimiento tácito.
\end{importantbox}

\subsection{Momento del cambio: oportunidades y riesgos}

A continuación se analiza el paso más arriesgado de la simplificación: cuándo puede eliminarse la ruta anterior. Para un fabricante de automóviles, esto no es un interruptor interno del grupo de investigación, sino una decisión que afecta directamente la seguridad de los usuarios y la calidad de la producción en serie.

Un fabricante de automóviles no puede lanzar algo solo porque la tecnología es atractiva. Para que el modelo sustituya de verdad a las reglas, hay que esperar a cierto punto: que el desempeño del modelo supere claramente al de las reglas, que los inconvenientes sean controlables, que los beneficios sean evidentes y que la responsabilidad de seguridad y la calidad de la producción en serie puedan sostenerse. Liu Xianming menciona el momento en que los probadores aplaudieron: esa es una señal organizacional de que el modelo superó a las reglas anteriores.

\begin{figure}[H]
\centering
\includegraphics[width=0.92\textwidth]{ai-transition-risk-balance.png}
\caption{Momento del cambio: oportunidades vs riesgos. El fabricante de automóviles debe esperar a que el modelo supere a las reglas, en lugar de lanzarlo solo porque la tecnología es atractiva.}
\end{figure}

\begin{warningbox}{El cambio de ruta tecnológica no puede imponerse por fe}
En un fabricante de automóviles, el cambio de ruta tecnológica debe responder por la seguridad, la calidad y los usuarios. Lo correcto no es «el modelo es atractivo, así que se lanza a toda la flota», sino lograr que el modelo supere al sistema anterior en los indicadores clave y ampliar el despliegue de forma gradual.
\end{warningbox}

\subsection{Resumen de la sección}

Esta transformación de Xiaopeng (小鹏) ocurre a la vez en dos planos, el técnico y el organizacional: en lo técnico se retiran las capas intermedias y en lo organizacional se acorta la cadena de información. Ambos sirven al mismo objetivo: que los datos y los problemas reales entren más rápido en la iteración del modelo.
